\documentclass{mynotes}

\begin{document}

\title{PGT theory derivation}
\author{Zhang Kai}
\date{\today}

\tableofcontents

\maketitle

\section{Fourier Diffraction Theorem (FDT)}
Under first Born approximation, the scattered field is given by the following integral:
\begin{equation}
	\psi_{\text{scat}}(\mathbf{r}) = \int G(\mathbf{r},\mathbf{r}') F(\mathbf{r}') \psi_{\text{inc}}(\mathbf{r}') \, \dd^3\mathbf{r}'.
\end{equation}

Here, \(G(\mathbf{r},\mathbf{r}')=\frac{\ee^{\ii k_m r}}{4\pi r}\) is the Green's function of the medium, where \(r = |\mathbf{r} - \mathbf{r}'|\).
The scattering potential \(F(\mathbf{r}')\) is defined as:
\begin{equation}
	F(\mathbf{r}') = k_m^2 \left( 1 - \frac{n^2(\mathbf{r}')}{n_m^2} \right).
\end{equation}

Under the out-going wave condition, and using the convolution theorem, the scattered field can be expressed in the Fourier domain as:
\begin{equation}
	\widetilde{\psi}_{\text{scat}}(\mathbf{k}) = \frac{\widetilde{F}(\mathbf{k}-\mathbf{k}_{\text{inc}})}{k^2 - k_m^2 - \ii 0^+}.
\end{equation}

\dots

We have the FDT (for the transmitted field) as follows:
\begin{equation}\label{eq:fdt}
	\widetilde{\psi}_{\text{scat}}(\mathbf{k}_{\perp},z)
	=\frac{\ee^{\ii \sqrt{k_m^2 - \abs{\mathbf{k}_{\perp}}^2} z}}
	{2\ii \sqrt{k_m^2 - \abs{\mathbf{k}_{\perp}}^2}}
	\widetilde{F}(\mathbf{k}_{\perp} - \mathbf{k}_{\text{inc}, \perp}, \sqrt{k_m^2 - \abs{\mathbf{k}_{\perp}}^2} - k_{\text{inc}, z}).
\end{equation}

In reality, the objective collecting the transmitted field has limited numerical aperture. The bandlimited-ness of the system is represented by the
pupil function

\begin{equation}
	P(\mathbf{k}_{\perp})
	=\begin{cases}
		1, & \abs{\mathbf{k}_{\perp}}<\text{NA}_{\text{det}}/\lambda \\
		0, & \text{else}
	\end{cases},
\end{equation}

the collected complex field in experiment is

\begin{equation}
	\widetilde{\psi}_{\text{scat}}(\mathbf{k}_{\perp},z)
	=P(\mathbf{k_{\perp}})
	\frac{\ee^{\ii \sqrt{k_m^2 - \abs{\mathbf{k}_{\perp}}^2} z}}
	{2\ii \sqrt{k_m^2 - \abs{\mathbf{k}_{\perp}}^2}}
	\widetilde{F}(\mathbf{k}_{\perp} - \mathbf{k}_{\text{inc}, \perp}, \sqrt{k_m^2 - \abs{\mathbf{k}_{\perp}}^2} - k_{\text{inc}, z}).
\end{equation}


\section{PGT theory derivation}

Take the Fourier transform with respect to \(z\) on both sides of Eq. \ref{eq:fdt}, we have\footnote{$2\pi$ factor in Fourier transfoming
a plane wave is omitted, because what matters is the ratio between terms.}:

\begin{equation}\label{eq:complextf}
	\widetilde{\psi}_{\text{scat}}(\mathbf{k})
	=P(\mathbf{k_{\perp}})
	\frac{\delta(k_z - \sqrt{k_m^2 - \abs{\mathbf{k}_{\perp}}^2})}{2 \ii k_z}
	\widetilde{F}(\mathbf{k}-\mathbf{k}_{\text{inc}}).
\end{equation}

Eq. \ref{eq:complextf} describes the linear relation between the scatter field and the scattering potential. To derive the complete PGT theory,
we have to take into consideration the unscattered field. In frequency domain,

\begin{equation}
	\widetilde{\psi}_{\text{tot}}(\mathbf{k})
	=P(\mathbf{k_{\perp}})\delta^{(3)}(\mathbf{k}-\mathbf{k}_{\text{inc}})
	+P(\mathbf{k_{\perp}})\frac{\delta(k_z - \sqrt{k_m^2 - \abs{\mathbf{k}_{\perp}}^2})}{2 \ii k_z} \widetilde{F}(\mathbf{k}-\mathbf{k}_{\text{inc}})
\end{equation}

The intensity stack captured in PGT experiment is related to the scatterd field by

\begin{equation}
	I(\mathbf{r})=\abs{\psi_{\text{tot}}(\mathbf{r})}^2
\end{equation}

Using the convolution theorem of Fourier transform, the spectrum of the intensity stack can be expressed as

\begin{equation}\label{eq:ispec}
	\widetilde{I}(\mathbf{k})=\widetilde{\psi}^*(-\mathbf{k})\ast\widetilde{\psi}(\mathbf{k}).
\end{equation}

In expanding Eq. \ref{eq:ispec}, we obtain three terms, the \text{i-i} term, the cross \text{i-s} term and the s-s term. In weak scattering regime, the s-s term is
ignored.

The two terms can contributes to the intensity stack significantly are

\begin{equation}
	\widetilde{I}_{\text{i-i}}
    =P^*(\mathbf{k_{\text{inc}, \perp}})P(\mathbf{k}_{\perp}+\mathbf{k}_{\text{inc}, \perp})\delta^{(3)}(\mathbf{k})
    =\abs{P(\mathbf{k_{\text{inc}, \perp}})}^2\delta^{(3)}(\mathbf{k}),
\end{equation}

and

\begin{equation}\label{eq:cross}
	\begin{aligned}
		\widetilde{I}_{\text{i-s}}
		= & \int\dd \mathbf{k'}^{3}\,\bigg[P^*(-\mathbf{k_{\perp}}')
		\delta^{(3)*}(-\mathbf{k}'-\mathbf{k}_{\text{inc}})
		P(\mathbf{k_{\perp}}-\mathbf{k_{\perp}}')
		\frac{\delta(k_z - k_z' - \sqrt{k_m^2 - \abs{\mathbf{k}_{\perp}-\mathbf{k}_{\perp}'}^2})}
		{2 \ii (k_z-k_z')} \widetilde{F}(\mathbf{k}-\mathbf{k}'-\mathbf{k}_{\text{inc}})              \\
		  & -P^*(-\mathbf{k_{\perp}}')
		\frac{\delta^{*}(-k_z' - \sqrt{k_m^2 - \abs{\mathbf{k}_{\perp}'}^2})}
		{2 \ii (-k_z')} \widetilde{F}^*(-\mathbf{k}'-\mathbf{k}_{\text{inc}})
		P(\mathbf{k_{\perp}}-\mathbf{k_{\perp}}')
		\delta^{(3)}(\mathbf{k}-\mathbf{k}'-\mathbf{k}_{\text{inc}})\bigg]                                 \\
		= & P^*(\mathbf{k_{\text{inc}, \perp}})P(\mathbf{k}_{\perp}+\mathbf{k}_{\text{inc}, \perp})
		\frac{\delta(k_z + k_{\text{inc}, z} - \sqrt{k_m^2 - \abs{\mathbf{k}_{\perp}+\mathbf{k}_{\text{inc}, \perp}}^2})}
		{2 \ii (k_z+k_{\text{inc},z})} \widetilde{F}(\mathbf{k})                                      \\
		  & -P(\mathbf{k_{\text{inc}, \perp}})P^*(\mathbf{k_{\text{inc}, \perp}}-\mathbf{k_{\perp}})
		\frac{\delta(k_{\text{inc}, z} - k_z - \sqrt{k_m^2 - \abs{\mathbf{k}_{\text{inc}, \perp}-\mathbf{k}_{\perp}}^2})}
		{2 \ii (k_{\text{inc}, z} - k_z)} \widetilde{F}^*(-\mathbf{k}).                                \\
	\end{aligned}
\end{equation}

Further break Eq. \ref{eq:cross} into the real part and imaginary part of $F$ to obtain:

\begin{equation}
    \begin{aligned}
		\widetilde{I}_{\text{i-s}}
		= & P^*(\mathbf{k_{\text{inc}, \perp}})P(\mathbf{k}_{\perp}+\mathbf{k}_{\text{inc}, \perp})
		\frac{\delta(k_z + k_{\text{inc}, z} - \sqrt{k_m^2 - \abs{\mathbf{k}_{\perp}+\mathbf{k}_{\text{inc}, \perp}}^2})}
		{2 \ii (k_z+k_{\text{inc},z})}
        \left[\widetilde{F}_{\text{re}}(\mathbf{k}) + \ii\widetilde{F}_{\text{im}}(\mathbf{k})\right]\\
		  & -P(\mathbf{k_{\text{inc}, \perp}})P^*(\mathbf{k_{\text{inc}, \perp}}-\mathbf{k_{\perp}})
		\frac{\delta(k_{\text{inc}, z} - k_z - \sqrt{k_m^2 - \abs{\mathbf{k}_{\text{inc}, \perp}-\mathbf{k}_{\perp}}^2})}
		{2 \ii (k_{\text{inc}, z} - k_z)}
        \left[\widetilde{F}_{\text{re}}(\mathbf{k}) - \ii\widetilde{F}_{\text{im}}(\mathbf{k})\right].                                \\
	\end{aligned}
\end{equation}

For real part of the scattering potential, 

\begin{equation}
    \begin{aligned}
		\widetilde{I}_{\text{i-s, re}}
		= -\frac{\ii}{2}
        \bigg[
        & P^*(\mathbf{k_{\text{inc}, \perp}})P(\mathbf{k}_{\perp}+\mathbf{k}_{\text{inc}, \perp})
		\frac{\delta(k_z + k_{\text{inc}, z} - \sqrt{k_m^2 - \abs{\mathbf{k}_{\perp}+\mathbf{k}_{\text{inc}, \perp}}^2})}
		{k_z+k_{\text{inc},z}}\\
		  & -P(\mathbf{k_{\text{inc}, \perp}})P^*(\mathbf{k_{\text{inc}, \perp}}-\mathbf{k_{\perp}})
		\frac{\delta(k_{\text{inc}, z} - k_z - \sqrt{k_m^2 - \abs{\mathbf{k}_{\text{inc}, \perp}-\mathbf{k}_{\perp}}^2})}
		{k_{\text{inc}, z} - k_z}
        \bigg] \widetilde{F}_{\text{re}}(\mathbf{k});\\
	\end{aligned}
\end{equation}
    
for the imaginary part of the scattering potential, 

\begin{equation}
    \begin{aligned}
		\widetilde{I}_{\text{i-s, im}}
		= \frac{1}{2}
        \bigg[
        & P^*(\mathbf{k_{\text{inc}, \perp}})P(\mathbf{k}_{\perp}+\mathbf{k}_{\text{inc}, \perp})
		\frac{\delta(k_z + k_{\text{inc}, z} - \sqrt{k_m^2 - \abs{\mathbf{k}_{\perp}+\mathbf{k}_{\text{inc}, \perp}}^2})}
		{k_z+k_{\text{inc},z}}\\
		  & +P(\mathbf{k_{\text{inc}, \perp}})P^*(\mathbf{k_{\text{inc}, \perp}}-\mathbf{k_{\perp}})
		\frac{\delta(k_{\text{inc}, z} - k_z - \sqrt{k_m^2 - \abs{\mathbf{k}_{\text{inc}, \perp}-\mathbf{k}_{\perp}}^2})}
		{k_{\text{inc}, z} - k_z}
        \bigg] \widetilde{F}_{\text{im}}(\mathbf{k}).\\
	\end{aligned}
\end{equation}

\end{document}
